\documentclass[10pt,a4paper]{amsart}
\usepackage[margin=19mm]{geometry}
\usepackage{amsmath,amssymb,mathtools}
\usepackage[scr=rsfs,cal=euler]{mathalfa}
\usepackage{xfrac}
\usepackage[unicode,hidelinks,bookmarks=false]{hyperref}
\newcommand{\ie}{i.e.\ }
\newcommand{\cf}{cf.\ }
\newcommand{\resp}{respectively\ }
\newcommand{\Subsection}{\subsection}
\providecommand{\nobreakdash}{\nobreak\hbox{-}}
\setlength{\emergencystretch}{2em}
\multlinegap0pt
\usepackage{bm}
\usepackage{bbm}
\usepackage{dsfont}
\usepackage{xspace}
\usepackage{cancel}
\usepackage{mathtools}
\usepackage{amsthm}
\makeatletter
\@ifundefined{proposition}{\newtheorem{proposition}{Proposition}}{}
\@ifundefined{theorem}{\newtheorem{theorem}{Theorem}}{}
\makeatother
\newcommand{\RR}{\mathbb R}
\newcommand{\R}{\mathbb R}
\newcommand{\argmax}{\mathop{\rm argmax}}
\newcommand{\bv}{\bar{v}}
\newcommand{\by}{\bar{y}}
\newcommand{\dist}{\mathtt{d}}
\newcommand{\inner}[2]{\langle#1,#2\rangle}
\newcommand{\norm}[1]{\textstyle{\Vert} #1 \textstyle{\Vert}}
\newcommand{\proj}{\mathtt{Proj}}
\title[Measuring dissimilarity between convex cones by means of max]{Measuring dissimilarity between convex cones by means of max-min angles}
\author{Welington de Oliveira, Valentina Sessa, David Sossa}
\date{Source: arXiv:2511.10483v2 (CC BY 4.0)}
\begin{document}
\maketitle

\noindent\emph{Source and licence.} Measuring dissimilarity between convex cones by means of max-min angles. Version arXiv:2511.10483v2 (CC BY 4.0), distributed under the Creative Commons Attribution 4.0 International licence, as recorded in the arXiv licence metadata for this exact version and shown on the abstract page https://arxiv.org/abs/2511.10483v2. Full rights notice: licenses/arxiv-2511.10483-CC-BY-4.0.txt.\par\smallskip

\noindent\textit{Editorial scope.} Counted unit: the Theorem whose source label is \texttt{None} in the article (source0.tex lines 608--616, source label \texttt{None}), delivered together with the article's own complete original proof of it (source0.tex lines 617--663, one \texttt{proof} environment of 47 lines). No mathematics is written, repaired or re-derived: statement and proof are the article's own text line for line. Required context retained uncounted: support (lines 444--446); prop FQ (lines 455--465); subspaces (lines 581--583). Every definition, display and label of those lines is retained unchanged. Omitted from this package and not redistributed: 1--443, 447--454, 466--580, 584--607, 664--1255. Nothing inside a retained excerpt is omitted or altered: every retained body is delivered exactly as the article's lines are written, so no omission receipt of any kind accompanies this package. Citations: the retained excerpts carry no citation command at all, so no bibliography entry of the article is transcribed and no reference is redistributed. Reference adaptation: none was needed and none was made; every reference inside the delivered unit resolves inside it, so no unresolved reference or citation is produced. Printed slips kept literal: no printed word, symbol, hypothesis or number of the counted statement or of its proof was altered. Layout and class: the article's own class is \texttt{siamart250211} with packages including lipsum, epstopdf, algorithmic, amscd, latexsym, graphicx, amsmath, amsfonts, float, amssymb, enumerate, multirow, floatflt, color; the wrapper is the lane's pinned \texttt{amsart} editorial wrapper at 10pt on A4, which restates the theorem-like environments the retained text needs and adds the article's own macro definitions that the retained text needs (\texttt{\textbackslash R}, \texttt{\textbackslash RR}, \texttt{\textbackslash argmax}, \texttt{\textbackslash bv}, \texttt{\textbackslash by}, \texttt{\textbackslash dist}, \texttt{\textbackslash inner}, \texttt{\textbackslash norm}, \texttt{\textbackslash proj}), with extra wrapper packages (bm, bbm, dsfont, xspace, cancel, mathtools). The delivered numbering is the unit's own. Assets: no retained span contains a figure environment, a graphics-inclusion command, a PDF-page inclusion, a table or a TikZ picture; no image file is copied into this package, the package declares no asset dependency and no image file is redistributed with it. Licence: version arXiv:2511.10483v2 of this article is distributed under the Creative Commons Attribution 4.0 International licence, as recorded in the arXiv licence metadata for this exact version and shown on the abstract page https://arxiv.org/abs/2511.10483v2. A full rights notice is delivered at licenses/arxiv-2511.10483-CC-BY-4.0.txt. Characters: every retained excerpt is pure ASCII, so the delivered engine, XeLaTeX with its default Latin Modern text font, reports no missing character.
\begin{equation}\label{support}
F_Q(u):=\max_{v\in Q\cap S_n}\langle u,v\rangle,
\end{equation}

\begin{proposition}\label{prop FQ}
Let $Q\in\mathcal C_n$, and let $u\in\R^n$. The following statements hold:
\begin{enumerate}
\item[(a)] The function $F_Q$ given in \eqref{support} is convex and its domain is $\R^n$.
\item[(b)]  $\displaystyle F_Q(u)=\frac{1+\norm{u}^2-\dist_{Q\cap S_n}^2(u)}{2}$.
\item[(c)] $\proj_{Q\cap S_n}(u)=\arg\max_{v \in Q\cap S_n} \inner{u}{v}$.
\item[(d)] $\partial F_Q(u) = \mathtt{conv}\, \proj_{Q\cap S_n}(u)$.
\item[(e)] $F_Q(u) = \inner{u}{v'}$ for all $v'\in \proj_{Q\cap S_n}(u)$.
\item[(f)] If $u \notin Q^\circ$, then $\partial F_Q(u)=\{\proj_Q(u)/\Vert \proj_Q(u)\Vert\}$. 
\end{enumerate}
\end{proposition}

\begin{equation}\label{subspaces}
P=U(\RR^p)\quad\text{and}\quad Q=V(\R^q),
\end{equation}

\begin{theorem}
Let $1\leq r\leq\infty$, and let $P=U(\R^p)$  and $Q=V(\R^q)$ be linear subspaces of $\R^n$ as in \eqref{subspaces}. Then,
\begin{align*}
&\displaystyle {\rm Dis}_r(P,Q)= 2\left\Vert\left(\frac{\sqrt{2}}{2}, \sin\left[\frac{\theta_{\max}(P,Q)}{2}\right]\right)\right\Vert_r,\quad\text{if }p \neq q,\\
&\displaystyle {\rm Dis}_r(P,Q)= 2\Vert(1,1)\Vert_r\sin\left[\frac{\theta_{\max}(P,Q)}{2}\right],\quad\text{if } p=q,
\end{align*}
where $\theta_{\max}(P,Q)\in[0,\pi/2]$ is the largest principal angle between $P$ and $Q$; that is, $\cos(\theta_{\max}(P,Q))$ is the smallest singular value of $V^\top U$.

\end{theorem}

\begin{proof}
Let $u\in\R^n$ be such that $\Vert u\Vert=1$. Let us compute the projection of $u$ onto $Q\cap S_n$:
\begin{eqnarray*}
\bv\in \proj_{Q\cap S_n}(u)&\Leftrightarrow&\bv\in\argmax_{v\in Q\cap S_n}\,\langle u,v\rangle\,\Leftrightarrow\, \bv=V\bar y\;\text{ with }\; \by\in\argmax_{y\in S_q}\,\langle y,V^\top u\rangle,\\
&\Leftrightarrow&\bv=\frac{VV^\top u}{\Vert V^\top u\Vert}\,\text{ if }\,V^\top u\neq 0; \,\text{otherwise}\;\bv\in Q\cap S_n.
\end{eqnarray*}
From Proposition\,\ref{prop FQ}(e) we know that $F_Q(u)=\langle u,\bv\rangle$ with $\bv\in \proj_{Q\cap S_n}(u)$. Therefore, from the above computation, we deduce that
\[F_Q(u)=\begin{cases}
\Vert V^\top u\Vert,&\text{if }u\notin Q^\perp\\
0,&\text{if }u\in Q^\perp
\end{cases}.\]
Hence, if $P\cap Q^\perp=\{0\}$,
\[\mathfrak{s}(P,Q)=\min_{u\in P\cap S_n}F_Q(u)=\min_{u\in P\cap S_n}\Vert V^\top u\Vert=\min_{x\in S_p}\Vert V^\top Ux\Vert=\sigma_{\min}(V^\top U),\]
where $\sigma_{\min}(V^\top U)$ denotes the smallest singular value of $V^\top U$. If $P\cap Q^\perp\neq\{0\}$, it is clear that $\mathfrak{s}(P,Q)=0$. Thus,

\[\mathfrak{s}(P,Q)=\begin{cases}
\sigma_{\min}(V^\top U),&\text{if }P\cap Q^\perp=\{0\},\\
0,&\text{if }P\cap Q^\perp\neq\{0\}.
\end{cases}\]
Recall that $\theta_{\max}(P,Q)=\arccos(\sigma_{\min}(V^\top U))$. Therefore, we deduce that
\begin{equation}\label{theta-arccos}
\Theta(P,Q)=\arccos(\mathfrak{s}(P,Q))=\begin{cases}
\theta_{\max}(P,Q),&\text{if }P\cap Q^\perp=\{0\},\\
\frac{\pi}{2},&\text{if }P\cap Q^\perp\neq\{0\}.
\end{cases}
\end{equation}
By interchanging the roles of $P$ and $Q$, and recalling that $\sigma_{\min}(V^\top U)=\sigma_{\min}(U^\top V)$, we also get that
\begin{equation}\label{theta-arccos2}
\Theta(Q,P)=\begin{cases}
\theta_{\max}(P,Q),&\text{if }P^\perp\cap Q=\{0\},\\
\frac{\pi}{2},&\text{if }P^\perp\cap Q\neq\{0\}.
\end{cases}
\end{equation}

Recall that $p=\mathtt{dim} P$ and $q=\mathtt{dim} Q$, and set $E:=V^\top U\in\R^{q\times p}$. Suppose that $p>q $. The Rank-Nullity Theorem says that $\mathtt{dim}({\rm Ker}(E))+\mathtt{rank}(E)=p$. Since $\mathtt{rank}(E)\leq q<p$, we deduce that $\mathtt{dim}({\rm Ker}(E))\geq 1$. It implies that there exists a nonzero $x\in\R^p$ such that $V^\top U x=0$. Thus, $Ux\in P\cap Q^\perp$, and from \eqref{theta-arccos} we deduce that $\Theta(P,Q)=\pi/2$. Now, to compute $\Theta(Q,P)$, suppose that $P^\perp\cap Q\neq\{0\}$. Then, there exists a nonzero $y\in\R^q$ such that $U^\top Vy=0$. Hence, $\mathtt{dim}({\rm Ker}(E^\top))\geq 1$. The Rank-Nullity Theorem for $E^\top$ says that $\mathtt{dim}({\rm Ker}(E^\top))+\mathtt{rank}(E)=q$. Then, $\mathtt{rank}(E)<q$ which means that $E$ is not of complete rank. Hence, $\sigma_{\min}(E)=0$, and then $\Theta(Q,P)=\theta_{\max}(P,Q)=\arccos(0)=\pi/2$. Therefore, by combining it with \eqref{theta-arccos2} we get that $\Theta(Q,P)=\theta_{\max}(P,Q)$ when either $P^\perp\cap Q\neq\{0\}$ or $P^\perp\cap Q=\{0\}$. Thus, 
\begin{equation}\label{dis}
\displaystyle {\rm Dis}_r(P,Q)= 2\left\Vert\left(\frac{\sqrt{2}}{2}, \sin\left[\frac{\theta_{\max}(P,Q)}{2}\right]\right)\right\Vert_r.
\end{equation}
By symmetry, the above result is also valid when $p<q$. Hence, \eqref{dis} holds whenever $p\neq q$.

Suppose that $p=q$. Then, analogous to the above reasoning, the matrix $E=V^\top U\in\R^{p\times p}$ is not of complete rank when $P\cap Q^\perp\neq\{0\}$. Thus, $\sigma_{\min}(E)=0$, and then $\Theta(P,Q)=\theta_{\max}(P,Q)=\arccos(0)=\pi/2$. By combining it with \eqref{theta-arccos} we get that $\Theta(P,Q)=\theta_{\max}(P,Q)$ in the both cases of \eqref{theta-arccos}. Analogously, we also have $\Theta(Q,P)=\theta_{\max}(P,Q)$. Hence, if $p=q$ then
\begin{eqnarray*}
\displaystyle {\rm Dis}_r(P,Q)&=& 2\left\Vert\left(\sin\left[\frac{\theta_{\max}(P,Q)}{2}\right], \sin\left[\frac{\theta_{\max}(P,Q)}{2}\right]\right)\right\Vert_r\\
&=&2\Vert(1,1)\Vert_r\sin\left[\frac{\theta_{\max}(P,Q)}{2}\right].
\end{eqnarray*}
The proof is complete.
\end{proof}

\end{document}
